% !TEX root = ../algorithms.tex

\section{Задача булевой выполнимости (SAT). Задачи $k$SAT и MAX-$k$SAT. Вероятностный и детерминированный приближённые алгоритмы для MAX-3SAT}

\subsection{Определения}

\begin{Definition}{}{}
	\emph{Литерал} --- переменная $x_i$ или её отрицание $\bar{x}_i$.
	\emph{Клауза} --- дизъюнкция литералов.
	\emph{КНФ} --- конъюнкция клауз.
	
	\textbf{SAT:} Дана КНФ с $n$ переменными и $m$ клаузами. Существует ли набор значений переменных, при котором она истинна? (\emph{NP-полная задача})
	
	\textbf{$k$SAT:} SAT для КНФ, где каждая клауза содержит ровно $k$ литералов.
	
	\textbf{MAX-$k$SAT:} Дана $k$-КНФ. Найти набор переменных, максимизирующий число истинных клауз. (\emph{NP-полная при $k \geq 2$})
\end{Definition}

\begin{Definition}{Задача MAX-3SAT}{}
	Дана булева формула в 3-КНФ:
	\[
	\varphi = C_1 \wedge C_2 \wedge \cdots \wedge C_m,
	\qquad
	C_j = (l_{j,1} \vee l_{j,2} \vee l_{j,3}),
	\]
	где каждый литерал $l_{j,i}$ равен $x_s$ или $\bar{x}_s$ для некоторого $s$.
	
	Требуется найти набор $x_1, \ldots, x_n \in \{0,1\}$, максимизирующий число выполненных клауз.
\end{Definition}

\subsection{Вероятностный алгоритм}

\textbf{Алгоритм.} Каждую переменную $x_i$ выбрать независимо и равновероятно из $\{0,1\}$.

\begin{Lemma}{}{}
	Пусть $X$ --- число выполненных клауз при независимом равновероятном выборе переменных. Тогда:
	\begin{enumerate}
		\item $\mathbb{E}[X] = \dfrac{7}{8}\,m$;
		\item $\mathbb{P}\!\left(X \geq \dfrac{7}{8}\,m\right) \geq \dfrac{1}{8m}$.
	\end{enumerate}
\end{Lemma}

\begin{proof}
	\textbf{Часть 1.}
	Клауза $C_j$ ложна тогда и только тогда, когда все три её литерала ложны, что происходит с вероятностью $\tfrac{1}{8}$. Значит, $\mathbb{P}(C_j \text{ выполнена}) = \tfrac{7}{8}$.
	По линейности ожидания: $\mathbb{E}[X] = m \cdot \tfrac{7}{8}$.
	
	\textbf{Часть 2.}
	Пусть $p_s = \mathbb{P}(X = s)$ и $p = \mathbb{P}\!\bigl(X \geq \lceil\tfrac{7m}{8}\rceil\bigr) = \sum_{s \geq \lceil 7m/8 \rceil} p_s$. Из $\mathbb{E}[X] = \tfrac{7m}{8}$:
	\begin{align*}
		\frac{7m}{8}
		&= \sum_{s=0}^{m} s\,p_s
		= \underbrace{\sum_{s < \lceil 7m/8 \rceil} s\,p_s}_{\leq\;(\lceil 7m/8 \rceil - 1)(1-p)}
		+ \underbrace{\sum_{s \geq \lceil 7m/8 \rceil} s\,p_s}_{\leq\; m\,p}
		\;\leq\; \bigl(\lceil\tfrac{7m}{8}\rceil - 1\bigr) + m\,p.
	\end{align*}
	Проверим, что $\lceil\tfrac{7m}{8}\rceil - 1 \leq \tfrac{7m}{8} - \tfrac{1}{8}$: пишем $7m = 8q + r$, $r \in \{0,\ldots,7\}$. При $r=0$: $q - 1 = \tfrac{7m}{8} - 1 \leq \tfrac{7m}{8} - \tfrac{1}{8}$. При $r \geq 1$: $q = \tfrac{7m-r}{8} \leq \tfrac{7m}{8} - \tfrac{1}{8}$.
	
	Подставляя: $\dfrac{7m}{8} \leq \dfrac{7m}{8} - \dfrac{1}{8} + m\,p$, откуда $p \geq \dfrac{1}{8m}$.
\end{proof}

\begin{Corollary}{}{}
	Для любой 3-КНФ существует набор переменных, выполняющий не менее $\tfrac{7}{8}m$ клауз.
\end{Corollary}

\subsubsection*{Вероятностный $\frac{8}{7}$-приближённый алгоритм}

Генерировать случайные наборы до тех пор, пока число выполненных клауз $\geq \tfrac{7}{8}m$; вернуть такой набор.

\begin{Theorem}{}{}
	Ожидаемое число итераций описанного алгоритма не превышает $8m$.
\end{Theorem}

\begin{proof}
	Каждая итерация успешна с вероятностью $p \geq \tfrac{1}{8m}$. Число итераций до первого успеха --- геометрическое: $\mathbb{E}[\text{итераций}] = \tfrac{1}{p} \leq 8m$.
\end{proof}

\subsection{Детерминированный алгоритм (дерандомизация методом условных ожиданий)}

\subsubsection*{Весовая функция}

Присваиваем переменным значения по одной: $x_1, x_2, \ldots, x_n$. После $k$-го шага частичное присваивание $x_1, \ldots, x_k$ уже зафиксировано.

\begin{Definition}{}{}
	Для каждой ещё не выполненной клаузы $C$ при частичном присваивании $x_1, \ldots, x_k$ её \emph{вес} равен
	\[
	w_C = \frac{1}{2^j},
	\]
	где $j$ --- число переменных клаузы $C$, которым ещё не присвоено значение:
	\[
	j = [i_1 > k] + [i_2 > k] + [i_3 > k].
	\]
	Здесь $[\cdot]$ --- скобки Айверсона. Суммарный вес:
	\[
	W^{(k)} = \sum_{\text{не выполнена } C} w_C.
	\]
\end{Definition}

\begin{Remark}{}{}
	$w_C$ равен вероятности того, что $C$ окажется невыполненной при случайном равновероятном завершении присваивания. Поэтому
	\[
	W^{(k)} = \mathbb{E}\bigl[\text{число невыполненных клауз} \mid x_1,\ldots,x_k\bigr].
	\]
	При $k = 0$: $j = 3$ для каждой клаузы, $w_C = \tfrac{1}{8}$, и $W^{(0)} = \tfrac{m}{8}$.
\end{Remark}

\subsubsection*{Жадный алгоритм}

На шаге $k+1$ присвоить $x_{k+1}$ то значение из $\{0,1\}$, при котором $W^{(k+1)}$ минимально.

\begin{Statement}{}{}
	При жадном выборе $W^{(k+1)} \leq W^{(k)}$.
\end{Statement}

\begin{proof}
	Пусть $W_0$ и $W_1$ --- суммарные веса невыполненных клауз при $x_{k+1} = 0$ и $x_{k+1} = 1$ соответственно. Покажем, что $W_0 + W_1 = 2W^{(k)}$, откуда $\min(W_0, W_1) \leq W^{(k)}$.
	
	Рассмотрим вклад каждой невыполненной клаузы $C$ (с текущим весом $w_C$) в $W_0 + W_1$:
	\begin{enumerate}
		\item \emph{$C$ не содержит $x_{k+1}$:} вес $w_C$ не меняется ни при $0$, ни при $1$. Вклад: $2w_C$.
		
		\item \emph{$C$ содержит позитивный литерал $x_{k+1}$:}
		при $x_{k+1} = 1$ клауза выполняется (вклад в $W_1$ равен $0$);
		при $x_{k+1} = 0$ число свободных переменных уменьшается на $1$, вес удваивается до $2w_C$.
		Вклад в $W_0 + W_1$: $2w_C + 0 = 2w_C$.
		
		\item \emph{$C$ содержит отрицательный литерал $\bar{x}_{k+1}$:} симметрично.
		При $x_{k+1} = 0$ клауза выполняется (вклад в $W_0$ равен $0$);
		при $x_{k+1} = 1$ вес удваивается.
		Вклад в $W_0 + W_1$: $0 + 2w_C = 2w_C$.
	\end{enumerate}
	
	Во всех случаях вклад равен $2w_C$, значит $W_0 + W_1 = 2W^{(k)}$, и $\min(W_0, W_1) \leq W^{(k)}$.
\end{proof}

\subsubsection*{Корректность}

\begin{Theorem}{}{}
	Жадный алгоритм дерандомизации за $n$ шагов находит набор значений, при котором выполнено не менее $\dfrac{7}{8}\,m$ клауз.
\end{Theorem}

\begin{proof}
	По предыдущему утверждению $W^{(k+1)} \leq W^{(k)}$ на каждом шаге, поэтому
	\[
	W^{(n)} \leq W^{(0)} = \frac{m}{8}.
	\]
	По завершении ($k = n$) все переменные присвоены: $j = 0$ для каждой клаузы, поэтому $w_C = 1$ для невыполненной и $w_C = 0$ для выполненной. Следовательно,
	\[
	\text{число невыполненных клауз} = W^{(n)} \leq \frac{m}{8},
	\]
	откуда число выполненных $\geq m - \dfrac{m}{8} = \dfrac{7m}{8}$.
\end{proof}

\begin{Remark}{}{}
	Время работы: на каждом из $n$ шагов перебирается два значения и вычисляется $W$ за $O(m)$. Итого $O(nm)$ --- полиномиально.
\end{Remark}

\begin{Remark}{}{}
	Хостад (2001) доказал, что $\bigl(\tfrac{8}{7} - \varepsilon\bigr)$-приближённый алгоритм для MAX-3SAT, работающий за полиномиальное время, влечёт $\mathrm{P} = \mathrm{NP}$.
\end{Remark}